% !TeX root = ../main.tex

\chapter{Experiments and Evaluation}\label{chapter:experiments-evaluation}

This chapter validates the control approach and aerodynamic performance of the X-wing platform through a series of experimental campaigns.
The experiments address four key objectives:
\begin{enumerate}
\item Identify the static thrust-to-RPM mapping for control and energy analysis (Munich, thrust stand).
\item Evaluate trajectory-tracking performance and agility of the INDI controller (Munich, flight arena).
\item Quantify aerodynamic forces and validate computational predictions (AIDA indoor flight tests).
\item Measure direct electrical power consumption under realistic flight conditions (outdoor tests).
\end{enumerate}

\section{Static Thrust Identification}\label{sec:thrust-identification}

The thrust-characterization procedure described in Section~\ref{sec:thrust-characterization} produced 975 measurements under the 6\,S 2\,P power configuration.
With that a quadratic model was fitted:
\begin{equation}
T = 5.15 \times 10^{-8}\,\text{RPM}^2 - 0.09 \times 10^{-3}\,\text{RPM},
\label{eq:thrust-rpm-fit}
\end{equation}
with coefficient of determination $R^2 = 0.997$ (Figure~\ref{fig:thrust_vs_rpm}).

The quadratic dependence confirms classical blade-element scaling $T \propto \Omega^2$ and validates the assumption used in the control allocation.
This mapping was used throughout all flight experiments to translate desired thrust into commanded motor speeds.

\begin{figure}[htbp]
  \centering
  \begin{tikzpicture}
  \begin{axis}[tumplot, xlabel={RPM [rev/min]}, ylabel={Thrust [N]}, width=0.8\textwidth]
      % Original scatter data points
      \addplot[only marks, mark=*, mark size=1pt, opacity=0.5, TUMBlue]
        table[x expr=1000*\thisrow{rpm_motor1}, y=force_z_N, col sep=comma]{data/thrust_map_raw_data.csv};
      \addlegendentry{Measured data}
      % Quadratic fit curve
      \addplot[thick, TUMAccentOrange, no marks]
        table[x expr=1000*\thisrow{rpm}, y=thrust_N_fitted, col sep=comma]{data/thrust_vs_rpm_quadratic_fit.csv};
      \addlegendentry{Quadratic fit ($R^2=0.997$)}
    \end{axis}
  \end{tikzpicture}
  \caption{Thrust vs RPM with zero-offset corrected quadratic fit. An offset of $-0.996$\,N was applied to ensure zero thrust at zero RPM.}
  \label{fig:thrust_vs_rpm}
\end{figure}

\begin{figure}[htbp]
  \centering
  \begin{tikzpicture}
  \begin{axis}[tumplot, xlabel={Command [\%]}, ylabel={RPM [rev/min]}, width=0.8\textwidth]
      % Original scatter data points
      \addplot[only marks, mark=*, mark size=1pt, opacity=0.5, TUMBlue]
        table[x=cmd, y expr=1000*\thisrow{rpm_motor1}, col sep=comma]{data/thrust_map_raw_data.csv};
      \addlegendentry{Measured data}
      % Linear fit line
      \addplot[thick, TUMAccentOrange, no marks]
        table[x=cmd, y expr=1000*\thisrow{rpm_fitted}, col sep=comma]{data/rpm_vs_cmd_linear_fit.csv};
      \addlegendentry{Linear fit ($R^2=0.991$)}
    \end{axis}
  \end{tikzpicture}
  \caption{RPM vs Command with linear fit: $\text{RPM} = 19260 \cdot \text{cmd} + 3650$.}
  \label{fig:rpm_vs_cmd}
\end{figure}

\begin{figure}[htbp]
  \centering
  \begin{tikzpicture}
  \begin{axis}[tumplot, xlabel={Command [-]}, ylabel={Thrust [N]}, width=0.8\textwidth]
      \addplot[only marks, mark=*, mark size=1pt, opacity=0.5, TUMBlue]
        table[x=cmd, y=force_z_N, col sep=comma]{data/thrust_vs_cmd_data.csv};
      \addlegendentry{Measured data}
    \end{axis}
  \end{tikzpicture}
  \caption{Thrust vs Command showing the direct relationship between motor command and thrust force.}
  \label{fig:thrust_vs_cmd}
\end{figure}

\section{Trajectory-Tracking Performance and Agility}\label{sec:tracking-performance}

To comprehensively evaluate the platform's control robustness and agility envelope, circular trajectory experiments were conducted at the flight arena in Munich.
These tests validate the INDI controller's ability to maintain stable tracking performance across a wide speed range while executing sustained coordinated turns with significant aerodynamic loading.

\subsection{Circular-Trajectory Tracking}\label{sec:circular-tracking}

Closed-loop tracking was evaluated on circular trajectories at cruise speeds between \SI{3}{\meter\per\second} and \SI{8}{\meter\per\second}.
Two reference circles were used: a \SI{2.0}{\meter} radius for lower speeds (\SI{3}{\meter\per\second} and \SI{5}{\meter\per\second}) and a \SI{2.3}{\meter} radius for higher speeds (\SI{6.5}{\meter\per\second} and \SI{8}{\meter\per\second}).
This resulted in centripetal accelerations ranging from \SI{4.5}{\meter\per\second\squared} at \SI{3}{\meter\per\second} to \SI{27.8}{\meter\per\second\squared} (approximately 2.8g) at \SI{8}{\meter\per\second}.

The root-mean-square (RMS) 3D position error increased from \SI{0.064}{\meter} at \SI{3}{\meter\per\second} to \SI{0.184}{\meter} at \SI{8}{\meter\per\second}, while maximum deviation remained below \SI{0.54}{\meter}.
These results demonstrate that the controller maintains sub-decimeter tracking accuracy despite the presence of aerodynamic lift and drag, confirming that explicit aerodynamic modeling is unnecessary within this speed envelope.

\subsection{Variable-Speed Agility Demonstration}\label{sec:circular-tracking-detailed}

\begin{figure}[htbp]
  \centering
  \begin{tikzpicture}
  \begin{axis}[
    tumplot,
    width=0.85\textwidth,
    height=0.85\textwidth,
    xlabel={X Position [\si{\meter}]},
    ylabel={Y Position [\si{\meter}]},
    axis equal,
    legend style={at={(0.5,0.5)}, anchor=center, font=\small},
    grid=major,
  ]
  % Reference circles
  \addplot[black, thick, dashed, domain=0:360, samples=100] 
    ({2.3*cos(x)}, {2.3*sin(x)});
  \addlegendentry{Reference (r=2.3m)}
  
  \addplot[gray, thick, dotted, domain=0:360, samples=100] 
    ({2.0*cos(x)}, {2.0*sin(x)});
  \addlegendentry{Reference (r=2.0m)}
  
  % Measured trajectories at different speeds
  \addplot[TUMBlue, thick, mark=none] 
    table[x=x, y expr=\thisrow{y}-0.3, col sep=comma]{data/circle/circle_trajectory_30mps.csv};
  \addlegendentry{\SI{3}{\meter\per\second}}
  
  \addplot[TUMAccentOrange, thick, mark=none] 
    table[x=x, y expr=\thisrow{y}-0.3, col sep=comma]{data/circle/circle_trajectory_50mps.csv};
  \addlegendentry{\SI{5}{\meter\per\second}}
  
  \addplot[TUMAccentGreen, thick, mark=none] 
    table[x=x, y expr=\thisrow{y}-0.3, col sep=comma]{data/circle/circle_trajectory_65mps.csv};
  \addlegendentry{\SI{6.5}{\meter\per\second}}
  
  \addplot[TUMAccentPink, thick, mark=none] 
    table[x=x, y expr=\thisrow{y}-0.3, col sep=comma]{data/circle/circle_trajectory_80mps.csv};
  \addlegendentry{\SI{8}{\meter\per\second}}
  \end{axis}
  \end{tikzpicture}
  \caption{Circular trajectory tracking at variable speeds (\SI{3}{\meter\per\second} to \SI{8}{\meter\per\second}). The \SI{3}{\meter\per\second} and \SI{5}{\meter\per\second} trajectories follow the \SI{2.0}{\meter} radius reference circle (gray dotted), while the \SI{6.5}{\meter\per\second} and \SI{8}{\meter\per\second} trajectories follow the \SI{2.3}{\meter} radius reference circle (black dashed). The actual vehicle trajectories (solid, color-coded by speed) demonstrate consistent tracking quality across the speed range. Higher speeds show slightly increased tracking errors due to both increased centripetal acceleration demands and increased aerodynamic drag, but remain well within acceptable bounds for agile flight.}
  \label{fig:variable_speed_trajectories}
\end{figure}

\paragraph{Tracking performance vs. speed.}
Figure~\ref{fig:tracking_error_vs_speed} quantifies how tracking accuracy varies with flight speed using two key metrics: the root-mean-square (RMS) error and the maximum error. The RMS error is defined as:
\begin{equation}
\text{RMS error} = \sqrt{\frac{1}{N}\sum_{i=1}^{N}e_i^2}
\end{equation}
where $e_i$ is the position error at time step $i$ and $N$ is the total number of samples. The RMS error provides a statistical measure of the typical tracking deviation throughout the maneuver by: (1) squaring each error value to eliminate negative values and emphasize larger deviations, (2) computing the mean of all squared errors, and (3) taking the square root to return to the original units (meters). Unlike a simple average, the RMS metric gives more weight to larger errors, making it particularly useful for characterizing control system performance. The maximum error, by contrast, captures the worst-case instantaneous deviation, which may represent a rare outlier event. Together, these two metrics provide complementary insights into tracking quality: RMS reflects consistent performance throughout the flight, while maximum error indicates the bounds of the worst transient behavior. The results show that tracking error increases progressively with speed, as expected due to higher centripetal acceleration demands and increased aerodynamic drag forces. However, even at \SI{8}{\meter\per\second}, both RMS and maximum errors remain well-bounded, demonstrating robust control performance across the entire agility envelope.

\begin{figure}[htbp]
  \centering
  \begin{tikzpicture}
  \begin{axis}[
    tumplot,
    ybar,
    bar width=12pt,
    width=0.85\textwidth,
    height=0.5\textwidth,
    xlabel={Flight Speed [\si{\meter\per\second}]},
    ylabel={Position Error [\si{\meter}]},
    legend pos=north west,
    ymin=0,
    xtick={3, 5, 6.5, 8},
    xticklabels={3, 5, 6.5, 8},
    enlarge x limits=0.15,
    legend style={font=\small},
  ]
  % RMS error (3D)
  \addplot[fill=TUMBlue, draw=TUMBlue] 
      table[x=speed, y=rms_error_3d, col sep=comma] 
      {data/circle/circle_tracking_errors.csv};
  \addlegendentry{RMS error (3D)}
  
  % RMS error (XY only)
  \addplot[fill=TUMAccentOrange, draw=TUMAccentOrange] 
      table[x=speed, y=rms_error_xy, col sep=comma] 
      {data/circle/circle_tracking_errors.csv};
  \addlegendentry{RMS error (XY only)}
  
  % Max error (3D)
  \addplot[fill=TUMAccentGreen, draw=TUMAccentGreen] 
      table[x=speed, y=max_error_3d, col sep=comma] 
      {data/circle/circle_tracking_errors.csv};
  \addlegendentry{Max error (3D)}
  
  % Max error (XY only)
  \addplot[fill=TUMAccentGray, draw=TUMAccentGray] 
      table[x=speed, y=max_error_xy, col sep=comma] 
      {data/circle/circle_tracking_errors.csv};
  \addlegendentry{Max error (XY only)}
  \end{axis}
  \end{tikzpicture}
  \caption{Tracking error as a function of flight speed during circular trajectory maneuvers. The 3D RMS error (blue) increases from \SI{0.064}{\meter} at \SI{3}{\meter\per\second} to \SI{0.184}{\meter} at \SI{8}{\meter\per\second}, while the XY-plane RMS error (orange) ranges from \SI{0.032}{\meter} to \SI{0.162}{\meter}. The 3D maximum error (green) increases from \SI{0.088}{\meter} to \SI{0.538}{\meter}, and the XY-plane maximum error (gray) ranges from \SI{0.056}{\meter} to \SI{0.454}{\meter}. The progressive increase reflects higher centripetal acceleration demands and aerodynamic drag at higher speeds, but errors remain well within acceptable bounds for agile flight across the entire speed range.}
  \label{fig:tracking_error_vs_speed}
\end{figure}

\paragraph{Bank angle and centripetal acceleration.}
Figure~\ref{fig:bank_angle_vs_speed} illustrates the aggressive maneuvering capability of the platform. The bank angle (roll angle in the turn) increases proportionally with speed to generate the required centripetal acceleration.

\begin{figure}[htbp]
  \centering
  \begin{tikzpicture}
  \begin{axis}[
    tumplot,
    width=0.85\textwidth,
    height=0.5\textwidth,
    xlabel={Cruise Speed [\si{\meter\per\second}]},
    ylabel={Bank Angle [\si{\degree}]},
    legend pos=north west,
    grid=major,
    ymin=0,
    ymax=90,
  ]
  % Measured maximum bank angle
  \addplot[TUMBlue, thick, mark=*, mark size=3pt] 
    table[x=speed, y=measured_max_bank, col sep=comma]{data/circle/circle_bank_angle_comparison.csv};
  \addlegendentry{Measured max bank angle}
  
  % Theoretical bank angle for coordinated turn
  \addplot[TUMAccentOrange, dashed, thick, mark=square*, mark size=3pt]
    table[x=speed, y=theoretical_bank, col sep=comma]{data/circle/circle_bank_angle_comparison.csv};
  \addlegendentry{Theoretical (coordinated turn)}
  \end{axis}
  \end{tikzpicture}
  \caption{Bank angle as a function of cruise speed during circular flight. The measured maximum bank angles (blue solid) increase from \SI{29.85}{\degree} at \SI{3}{\meter\per\second} to \SI{76.77}{\degree} at \SI{8}{\meter\per\second}, closely tracking the theoretical predictions (orange dashed) for coordinated turns where $\tan(\phi) = a_c/g$. The close agreement validates the coordinated turn control strategy, with the platform achieving bank angles within \SI{2}{\degree}--\SI{10}{\degree} of the theoretical values across the entire speed range. At the highest speed (\SI{8}{\meter\per\second}), the platform sustains a bank angle of nearly \SI{77}{\degree}, corresponding to a centripetal acceleration of \SI{27.83}{\meter\per\second\squared} (approximately 2.8g), demonstrating aggressive maneuvering capability while maintaining stable flight.}
  \label{fig:bank_angle_vs_speed}
\end{figure}

\begin{table}[htbp]
\centering
\caption{Summary of tracking performance at different cruise speeds.}
\label{tab:variable_speed_metrics}
\begin{tabular}{lccccc}
\hline
\textbf{Speed} & \textbf{Centripetal} & \textbf{Bank} & \textbf{RMS} & \textbf{Max} & \textbf{Max Motor} \\
\textbf{[\si{\meter\per\second}]} & \textbf{Accel. [\si{\meter\per\second\squared}]} & \textbf{Angle [\si{\degree}]} & \textbf{Error [\si{\meter}]} & \textbf{Error [\si{\meter}]} & \textbf{Speed [\si{rpm}]} \\
\hline
3.0 & 4.5 & 29.8 & 0.064 & 0.088 & 1319 \\
5.0 & 12.5 & 57.3 & 0.238 & 0.310 & 1588 \\
6.5 & 18.4 & 70.4 & 0.155 & 0.386 & 1857 \\
8.0 & 27.8 & 76.8 & 0.184 & 0.538 & 1961 \\
\hline
\end{tabular}
\end{table}

\paragraph{Detailed attitude response at maximum speed.}
Figure~\ref{fig:attitude_8mps} presents the attitude time series for one complete circle at \SI{8}{\meter\per\second}, demonstrating the quality of the coordinated turn execution.

\begin{figure}[htbp]
  \centering
  \begin{tikzpicture}
  \begin{axis}[
    tumplot,
    width=0.95\textwidth,
    height=0.45\textwidth,
    xlabel={Time [\si{\second}]},
    ylabel={Attitude Angle [\si{\degree}]},
    legend pos=north west,
  ]
  % Roll angle (bank angle in turn) - measured and reference
  \addplot[TUMBlue, thick, mark=none] 
    table[x=time, y=roll_measured, col sep=comma]{data/circle/circle_attitude_8mps.csv};
  \addlegendentry{Roll (measured)}
  
  \addplot[TUMBlue, thick, dashed, mark=none] 
    table[x=time, y=roll_reference, col sep=comma]{data/circle/circle_attitude_8mps.csv};
  \addlegendentry{Roll (reference)}
  
  % Pitch angle
  \addplot[TUMAccentOrange, thick, mark=none] 
    table[x=time, y=pitch_measured, col sep=comma]{data/circle/circle_attitude_8mps.csv};
  \addlegendentry{Pitch (measured)}
  
  \addplot[TUMAccentOrange, thick, dashed, mark=none] 
    table[x=time, y=pitch_reference, col sep=comma]{data/circle/circle_attitude_8mps.csv};
  \addlegendentry{Pitch (reference)}
  
  % Yaw angle (should rotate 360 deg)
  \addplot[TUMAccentGreen, thick, mark=none] 
    table[x=time, y=yaw_measured, col sep=comma]{data/circle/circle_attitude_8mps.csv};
  \addlegendentry{Yaw (measured)}
  
  \addplot[TUMAccentGreen, thick, dashed, mark=none] 
    table[x=time, y=yaw_reference, col sep=comma]{data/circle/circle_attitude_8mps.csv};
  \addlegendentry{Yaw (reference)}
  \end{axis}
  \end{tikzpicture}
  \caption{Attitude time series for one complete circle at \SI{8}{\meter\per\second}. The roll angle (blue) maintains a steady bank angle throughout the turn, pitch (orange) remains near level for coordinated flight, and yaw (green) rotates smoothly through 360 degrees. Both measured (solid) and reference (dashed) angles are shown.}
  \label{fig:attitude_8mps}
\end{figure}

\subsection{Agility and Disturbance Rejection}\label{sec:agility}

The vehicle remained controllable throughout the tested acceleration range, sustaining coordinated turns up to approximately 2.8\,g centripetal acceleration without attitude saturation or oscillation.
No divergence or limit-cycle behavior was observed, indicating that the incremental feedback successfully compensates for the additional aerodynamic moments induced by the wings.

\paragraph{Summary of agility demonstration.}
The variable-speed circular trajectory tests demonstrate that the X-wing platform achieves robust tracking performance across a wide speed range from \SI{3}{\meter\per\second} to \SI{8}{\meter\per\second}. Key findings include:

\begin{itemize}
  \item \textbf{Consistent tracking:} Position tracking errors remain bounded across all speeds, with RMS errors below \SI{0.24}{\meter} even at maximum speed
  \item \textbf{Coordinated turn behavior:} Bank angles closely follow theoretical predictions, validating the coordinated-turn control law implementation
  \item \textbf{Smooth attitude control:} Attitude time series show smooth, periodic behavior without oscillations or control saturation
  \item \textbf{Wide agility envelope:} The platform can sustain centripetal accelerations up to \SI{27.8}{\meter\per\second\squared} (approximately 2.8g) while maintaining stable flight
  \item \textbf{Control robustness:} The INDI-based controller successfully rejects aerodynamic disturbances from the wings across all flight speeds without requiring speed-dependent gain scheduling
\end{itemize}

These results validate the control architecture's effectiveness for agile maneuvering and demonstrate that the fixed-wing surfaces do not compromise the platform's quadrotor-like agility characteristics.

\section{Aerodynamic Coefficient Identification}\label{sec:aerodynamic_forces_quantification}

To validate computational aerodynamic predictions and quantify the actual forces experienced during flight, experimental coefficient identification was performed at the AIDA flight arena using force balance analysis from flight test data.
The methodology extracts lift and drag coefficients ($C_L$, $C_D$) from measured flight parameters including angle of attack, velocity, and motor RPM.

The aerodynamic coefficient analysis was conducted for both wing designs to quantify the effect of airfoil geometry on lift and drag.
The NACA~0015 configuration was used for all circular and agility tests, while both NACA~0015 and AG25 configurations were evaluated for coefficient identification.
The subsequent outdoor flight experiments employed the AG25 variant, whose superior lift-to-drag characteristics were confirmed by this analysis.

\paragraph{Rotor unloading measurements.}
Table~\ref{tab:rotor_unloading} summarizes the measured rotor speeds during straight-line flight tests at different cruise speeds for the NACA~0015 platform (Wing~1). All measurements use a constant hover RPM of \SI{1200}{rpm} as the baseline for comparison.

\begin{table}[htbp]
\centering
\caption{Rotor unloading during forward flight for NACA~0015 platform (Wing~1). Hover RPM: \SI{1200}{rpm}.}
\label{tab:rotor_unloading}
\begin{tabular}{cccccc}
\hline
\textbf{Flight Speed} & \textbf{Pitch Angle} & \textbf{Cruise RPM} & \textbf{RPM Ratio} & \textbf{Power Ratio} & \textbf{Power Savings} \\
\textbf{[\si{\meter\per\second}]} & \textbf{[\si{\degree}]} & \textbf{[\si{rpm}]} & \textbf{[-]} & \textbf{[-]} & \textbf{[\%]} \\
\hline
9.1  & 37.2 & 1130 & 0.942 & 0.836 & 16.4 \\
10.7 & 28.5 & 1110 & 0.925 & 0.792 & 20.8 \\
11.4 & 22.3 & 1050 & 0.875 & 0.670 & 33.0 \\
13.6 & 16.3 &  997 & 0.831 & 0.574 & 42.6 \\
\hline
\end{tabular}
\end{table}

\paragraph{Methodology.}
For steady-level flight, the force balance equations yield:
\begin{align}
L &= W - T \sin(\alpha) \label{eq:lift_balance} \\
D &= T \cos(\alpha) \label{eq:drag_balance}
\end{align}
where $W$ is the gravitational force (vehicle weight), $T$ is the total thrust (estimated from motor RPM using the thrust map from Section~\ref{sec:thrust-identification}), and $\alpha$ is the angle of attack measured from the horizontal plane. This angle is obtained by subtracting the quadrotor pitch angle from \SI{90}{\degree} ($\alpha = \SI{90}{\degree} - \theta_{\text{pitch}}$), which aligns the angle convention with standard fixed-wing aircraft notation and facilitates direct comparison with the XFLR5 airfoil analysis. The aerodynamic coefficients are then computed as:
\begin{align}
C_L &= \frac{L}{\frac{1}{2}\rho V^2 S} \label{eq:CL_definition} \\
C_D &= \frac{D}{\frac{1}{2}\rho V^2 S} \label{eq:CD_definition}
\end{align}
where $\rho = \SI{1.225}{\kilogram\per\meter\cubed}$ is air density, $V$ is flight speed, and $S$ is the wing reference area.

\paragraph{Experimental results.}
Flight tests were conducted with both wing configurations at various speeds. NACA~0015 tests covered speeds from \SI{9.14}{\meter\per\second} to \SI{13.6}{\meter\per\second} with pitch angles from \SI{37.2}{\degree} to \SI{16.3}{\degree}, while AG25 tests were performed at \SI{9.1}{\meter\per\second} and \SI{12.5}{\meter\per\second} with pitch angles from \SI{15.3}{\degree} to \SI{3.5}{\degree}. The extracted coefficients are integrated into Figures~\ref{fig:cl_comparison}--\ref{fig:ld_comparison} as markers (NACA~0015: green circles, AG25: pink squares), showing:

\begin{itemize}
  \item \textbf{Lift coefficients}: Experimental $C_L$ values show excellent agreement with XFLR5 predictions across the operational angle-of-attack range for both wing configurations, validating the computational model for lift estimation and confirming accurate force balance extraction. The AG25 achieves higher lift at significantly lower angles of attack compared to NACA~0015.
  \item \textbf{Drag coefficients}: Measured $C_D$ values are consistently higher than the XFLR5 wing-only simulations for both configurations, as expected since the simulation captures only the wing drag while the experimental data includes additional system-level drag sources such as fuselage parasite drag, propeller-induced drag, propeller-wing interference effects, and landing gear.
  \item \textbf{L/D ratios}: Experimental efficiency is lower than isolated wing performance due to these additional drag contributions, but the trends validate the force balance approach and demonstrate measurable aerodynamic contribution from both wing designs. The AG25 experimental data confirms superior efficiency at lower angles of attack compared to NACA~0015.
\end{itemize}

\paragraph{Discussion and sensitivity.}
The primary sources of uncertainty in the experimental coefficient identification include:
\begin{itemize}
  \item \textbf{Angle of attack estimation}: Pitch angle from IMU assumes level flight; any vertical velocity introduces error in effective angle of attack
  \item \textbf{Thrust estimation}: Motor RPM-to-thrust mapping has $\pm$5\% uncertainty from the thrust stand calibration
  \item \textbf{Steady-state assumption}: Transient accelerations violate the force balance equations, requiring careful selection of quasi-steady flight segments
\end{itemize}

Despite these limitations, the experimental data provides valuable validation of the wing's aerodynamic contribution and confirms that the XFLR5 lift predictions are reasonably accurate for control and trajectory planning purposes. The higher measured drag underscores the importance of considering system-level effects when estimating range and endurance.

\subsection{Computational Aerodynamic Characterization}\label{sec:airfoil_comparison}

To evaluate the potential for improved efficiency through airfoil optimization and validate the experimental measurements, we analyzed two wing configurations using XFLR5: the baseline NACA~0015 symmetric airfoil and an improved AG25 profile.
The analysis covers Reynolds numbers of $Re=1\times10^5$, $Re=2\times10^5$, and $Re=5\times10^5$, representing the operational flight regime from low to moderate speeds.
The primary operational condition at \SI{10}{\meter\per\second} cruise speed corresponds to $Re=2\times10^5$.
All simulations use NCrit=5 to account for the free-stream turbulence typical of indoor and outdoor flight environments.

Figures~\ref{fig:cl_comparison}--\ref{fig:ld_comparison} present comprehensive aerodynamic comparisons at two Reynolds numbers ($Re=1\times10^5$ and $Re=5\times10^5$).
The results demonstrate excellent aerodynamic performance at operational Reynolds numbers, with the AG25 airfoil achieving lift-to-drag ratios exceeding 50 at low Reynolds numbers and above 60 at the cruise condition.
The experimental data points, shown as markers in the figures, demonstrate good agreement with the theoretical predictions, particularly in the operational flight regime, validating both the XFLR5 model and the force balance extraction methodology.

\begin{figure}[htbp]
\centering
\input{figures/pgfplots_styles.tex}
\begin{tikzpicture}
\begin{axis}[
  tumplot,
  width=0.85\textwidth,
  height=0.5\textwidth,
  xlabel={Angle of Attack, $\alpha$ [\si{\degree}]},
  ylabel={Lift Coefficient, $C_L$ [-]},
  legend style={
    at={(0.5,1.05)},
    anchor=south,
    legend columns=2
  },
  grid=major,
]
% XFLR5 predictions
\addplot[tumblue, thick] table[x=alpha, y=CL, col sep=comma] {data/airfoil_analysis/CL_vs_alpha_NACA0015_Re0.100.csv};
\addlegendentry{NACA 0015, $Re=1\times10^5$ (XFLR5)}
\addplot[tumblue, thick, dashed] table[x=alpha, y=CL, col sep=comma] {data/airfoil_analysis/CL_vs_alpha_NACA0015_Re0.500.csv};
\addlegendentry{NACA 0015, $Re=5\times10^5$ (XFLR5)}
\addplot[tumorange, thick] table[x=alpha, y=CL, col sep=comma] {data/airfoil_analysis/CL_vs_alpha_AG25_Re0.100.csv};
\addlegendentry{AG25, $Re=1\times10^5$ (XFLR5)}
\addplot[tumorange, thick, dashed] table[x=alpha, y=CL, col sep=comma] {data/airfoil_analysis/CL_vs_alpha_AG25_Re0.500.csv};
\addlegendentry{AG25, $Re=5\times10^5$ (XFLR5)}
% Experimental data
\addplot[only marks, mark=*, mark size=4pt, tumblue, mark options={solid}] 
  table[x=alpha_deg, y=CL, col sep=comma] {data/aerodynamic_analysis/NACA0015_experimental_aerodynamic_coefficients.csv};
\addlegendentry{NACA 0015 (Experimental)}
\addplot[only marks, mark=square*, mark size=4pt, tumorange, mark options={solid}] 
  table[x=alpha_deg, y=CL, col sep=comma] {data/aerodynamic_analysis/AG25_experimental_aerodynamic_coefficients.csv};
\addlegendentry{AG25 (Experimental)}
\end{axis}
\end{tikzpicture}
\caption{Lift coefficient comparison between NACA~0015 and AG25 airfoils at two Reynolds numbers ($Re=1\times10^5$ and $Re=5\times10^5$), computed using XFLR5 with NCrit=5. Experimental data points from flight tests show excellent agreement with XFLR5 predictions for both wing configurations: NACA~0015 (green circles) and AG25 (pink squares), validating the simulation model for lift estimation. At $Re=1\times10^5$, the AG25 achieves maximum \(C_L \approx 1.20\) at \SI{10}{\degree}, while the NACA~0015 reaches \(C_L \approx 0.97\) at \SI{12}{\degree}.}
\label{fig:cl_comparison}
\end{figure}

\begin{figure}[htbp]
\centering
\input{figures/pgfplots_styles.tex}
\begin{tikzpicture}
\begin{axis}[
  tumplot,
  width=0.85\textwidth,
  height=0.5\textwidth,
  xlabel={Angle of Attack, $\alpha$ [\si{\degree}]},
  ylabel={Drag Coefficient, $C_D$ [-]},
  legend style={
    at={(0.5,1.05)},
    anchor=south,
    legend columns=2
  },
  grid=major,
]
% XFLR5 predictions
\addplot[tumblue, thick] table[x=alpha, y=CD, col sep=comma] {data/airfoil_analysis/CD_vs_alpha_NACA0015_Re0.100.csv};
\addlegendentry{NACA 0015, $Re=1\times10^5$ (XFLR5)}
\addplot[tumblue, thick, dashed] table[x=alpha, y=CD, col sep=comma] {data/airfoil_analysis/CD_vs_alpha_NACA0015_Re0.500.csv};
\addlegendentry{NACA 0015, $Re=5\times10^5$ (XFLR5)}
\addplot[tumorange, thick] table[x=alpha, y=CD, col sep=comma] {data/airfoil_analysis/CD_vs_alpha_AG25_Re0.100.csv};
\addlegendentry{AG25, $Re=1\times10^5$ (XFLR5)}
\addplot[tumorange, thick, dashed] table[x=alpha, y=CD, col sep=comma] {data/airfoil_analysis/CD_vs_alpha_AG25_Re0.500.csv};
\addlegendentry{AG25, $Re=5\times10^5$ (XFLR5)}
% Experimental data
\addplot[only marks, mark=*, mark size=4pt, tumblue, mark options={solid}] 
  table[x=alpha_deg, y=CD, col sep=comma] {data/aerodynamic_analysis/NACA0015_experimental_aerodynamic_coefficients.csv};
\addlegendentry{NACA 0015 (Experimental)}
\addplot[only marks, mark=square*, mark size=4pt, tumorange, mark options={solid}] 
  table[x=alpha_deg, y=CD, col sep=comma] {data/aerodynamic_analysis/AG25_experimental_aerodynamic_coefficients.csv};
\addlegendentry{AG25 (Experimental)}
\end{axis}
\end{tikzpicture}
\caption{Drag coefficient comparison between NACA~0015 and AG25 airfoils at two Reynolds numbers. Experimental data from flight tests (green markers) shows higher drag than XFLR5 wing-only simulations, as expected since the experimental values include additional system-level drag sources not present in the simulation: fuselage parasite drag, propeller-induced drag, propeller-wing interference, and landing gear. The XFLR5 predictions represent wing drag only. At $Re=1\times10^5$, the AG25 wing achieves minimum \(C_D \approx 0.011\) near \(\alpha = \SI{-1}{\degree}\), significantly lower than the NACA~0015's \(C_D \approx 0.015\) at \(\alpha = \SI{0}{\degree}\). The AG25 demonstrates superior low-drag characteristics across the operational angle of attack range.}
\label{fig:cd_comparison}
\end{figure}

\begin{figure}[htbp]
\centering
\input{figures/pgfplots_styles.tex}
\begin{tikzpicture}
\begin{axis}[
  tumplot,
  width=0.85\textwidth,
  height=0.5\textwidth,
  xlabel={Drag Coefficient, $C_D$ [-]},
  ylabel={Lift Coefficient, $C_L$ [-]},
  legend style={
    at={(0.5,1.05)},
    anchor=south,
    legend columns=2
  }
]
% XFLR5 predictions
\addplot[tumblue, thick] table[x=CD, y=CL, col sep=comma] {data/airfoil_analysis/drag_polar_NACA0015_Re0.100.csv};
\addlegendentry{NACA 0015, $Re=1\times10^5$ (XFLR5)}
\addplot[tumblue, thick, dashed] table[x=CD, y=CL, col sep=comma] {data/airfoil_analysis/drag_polar_NACA0015_Re0.500.csv};
\addlegendentry{NACA 0015, $Re=5\times10^5$ (XFLR5)}
\addplot[tumorange, thick] table[x=CD, y=CL, col sep=comma] {data/airfoil_analysis/drag_polar_AG25_Re0.100.csv};
\addlegendentry{AG25, $Re=1\times10^5$ (XFLR5)}
\addplot[tumorange, thick, dashed] table[x=CD, y=CL, col sep=comma] {data/airfoil_analysis/drag_polar_AG25_Re0.500.csv};
\addlegendentry{AG25, $Re=5\times10^5$ (XFLR5)}
% Experimental data
\addplot[only marks, mark=*, mark size=4pt, tumblue, mark options={solid}] 
  table[x=CD, y=CL, col sep=comma] {data/aerodynamic_analysis/NACA0015_experimental_aerodynamic_coefficients.csv};
\addlegendentry{NACA 0015 (Experimental)}
\addplot[only marks, mark=square*, mark size=4pt, tumorange, mark options={solid}] 
  table[x=CD, y=CL, col sep=comma] {data/aerodynamic_analysis/AG25_experimental_aerodynamic_coefficients.csv};
\addlegendentry{AG25 (Experimental)}
\end{axis}
\end{tikzpicture}
\caption{Drag polar comparison showing the relationship between lift and drag coefficients. Curves shifted toward the left and top indicate better aerodynamic efficiency (higher lift for given drag). Experimental data (green markers) achieves lift coefficients matching XFLR5 predictions but with higher drag coefficients, as expected since the experimental values include complete system drag (fuselage, propellers, landing gear, and propeller-wing interference) while XFLR5 simulates wing drag only. At both Reynolds numbers, the AG25 demonstrates dramatically superior wing efficiency with the drag polar shifted significantly leftward compared to NACA~0015, achieving higher lift coefficients at substantially lower drag penalties.}
\label{fig:drag_polar}
\end{figure}

\begin{figure}[htbp]
\centering
\input{figures/pgfplots_styles.tex}
\begin{tikzpicture}
\begin{axis}[
  tumplot,
  width=0.85\textwidth,
  height=0.5\textwidth,
  xlabel={Angle of Attack, $\alpha$ [\si{\degree}]},
  ylabel={Lift-to-Drag Ratio, $L/D$ [-]},
  legend style={
    at={(0.5,1.05)},
    anchor=south,
    legend columns=2
  },
  grid=major,
]
% XFLR5 predictions
\addplot[tumblue, thick] table[x=alpha, y=LD, col sep=comma] {data/airfoil_analysis/LD_vs_alpha_NACA0015_Re0.100.csv};
\addlegendentry{NACA 0015, $Re=1\times10^5$ (XFLR5)}
\addplot[tumblue, thick, dashed] table[x=alpha, y=LD, col sep=comma] {data/airfoil_analysis/LD_vs_alpha_NACA0015_Re0.500.csv};
\addlegendentry{NACA 0015, $Re=5\times10^5$ (XFLR5)}
\addplot[tumorange, thick] table[x=alpha, y=LD, col sep=comma] {data/airfoil_analysis/LD_vs_alpha_AG25_Re0.100.csv};
\addlegendentry{AG25, $Re=1\times10^5$ (XFLR5)}
\addplot[tumorange, thick, dashed] table[x=alpha, y=LD, col sep=comma] {data/airfoil_analysis/LD_vs_alpha_AG25_Re0.500.csv};
\addlegendentry{AG25, $Re=5\times10^5$ (XFLR5)}
% Experimental data
\addplot[only marks, mark=*, mark size=4pt, tumblue, mark options={solid}] 
  table[x=alpha_deg, y=L_D_ratio, col sep=comma] {data/aerodynamic_analysis/NACA0015_experimental_aerodynamic_coefficients.csv};
\addlegendentry{NACA 0015 (Experimental)}
\addplot[only marks, mark=square*, mark size=4pt, tumorange, mark options={solid}] 
  table[x=alpha_deg, y=L_D_ratio, col sep=comma] {data/aerodynamic_analysis/AG25_experimental_aerodynamic_coefficients.csv};
\addlegendentry{AG25 (Experimental)}
\end{axis}
\end{tikzpicture}
\caption{Lift-to-drag ratio comparison highlighting the aerodynamic efficiency of both airfoils. Experimental data (green markers) shows lower L/D ratios than XFLR5 wing-only predictions due to additional system drag sources (fuselage, propellers, landing gear) not included in the wing simulation. The XFLR5 curves represent wing performance only. At $Re=1\times10^5$, the AG25 wing achieves a remarkable maximum L/D of 53.0 at $\alpha \approx \SI{5}{\degree}$, with L/D ratios exceeding 50 between \SI{4}{\degree} and \SI{6}{\degree}—substantially higher than the NACA~0015 wing's maximum L/D of 37.1 at \SI{6}{\degree}. The experimental measurements at higher angles of attack (\SI{15}{\degree}--\SI{40}{\degree}) reflect the quadrotor flight regime where aerodynamic efficiency is secondary to hover capability.}
\label{fig:ld_comparison}
\end{figure}

\paragraph{Key aerodynamic findings.}
Table~\ref{tab:airfoil_metrics} summarizes the key performance metrics for both airfoils at $Re=2\times10^5$, which corresponds to the typical cruise speed of \SI{10}{\meter\per\second}.

\begin{table}[htbp]
\centering
\caption{Aerodynamic performance metrics comparison at $Re=2\times10^5$ (NCrit=5).}
\label{tab:airfoil_metrics}
\begin{tabular}{lcc}
\hline
\textbf{Metric} & \textbf{NACA 0015} & \textbf{AG25} \\
\hline
Maximum $L/D$ & 46.4 at \SI{7}{\degree} & 66.4 at \SI{5}{\degree} \\
$C_L$ at max $L/D$ & 0.81 & 0.83 \\
$C_D$ at max $L/D$ & 0.018 & 0.013 \\
Minimum $C_D$ & 0.010 at \SI{0}{\degree} & 0.008 at \SI{0}{\degree} \\
Maximum $C_L$ & 1.05 at \SI{13}{\degree} & 1.27 at \SI{11}{\degree} \\
$L/D$ at \SI{4}{\degree} cruise & 34.3 & 66.2 \\
$L/D$ at \SI{6}{\degree} cruise & 43.9 & 64.4 \\
\hline
\end{tabular}
\end{table}

The AG25 airfoil demonstrates a 43\% improvement in maximum lift-to-drag ratio compared to the NACA~0015 (66.4 vs 46.4), achieved at a lower angle of attack (\SI{5}{\degree} vs \SI{7}{\degree}).
At typical cruise conditions between \SI{4}{\degree} and \SI{6}{\degree}, the AG25 maintains L/D ratios of 64--66.
The AG25 also achieves 20\% lower minimum drag coefficient (0.008 vs 0.010) and 21\% higher maximum lift coefficient (1.27 vs 1.05).

\section{Outdoor Flight Validation}\label{sec:outdoor-flight}

The outdoor power-efficiency tests were performed using the improved AG25 $\pm\SI{20}{\degree}$ configuration, chosen based on the aerodynamic-coefficient analysis that demonstrated its superior efficiency compared to the baseline NACA~0015 $\pm\SI{45}{\degree}$ design.

To validate the platform's performance in realistic operational conditions beyond the controlled indoor environment, outdoor flight tests were conducted with electrical power logging and GPS-based velocity measurement. These tests provide direct evidence of aerodynamic efficiency by measuring actual power consumption during cruise flight.

\subsection{Test Configuration and Methodology}\label{sec:outdoor-methodology}

The outdoor tests were performed using the AG25-equipped platform (Wing~2, mass \SI{2.8}{\kg}) in calm open-field conditions with manual flight control and onboard electrical power logging. Environmental conditions were: wind speed < \SI{2}{\meter\per\second}, temperature \SI{20}{\celsius}, barometric pressure \SI{1013}{\hecto\pascal}.

Flight data was recorded at \SI{1}{\kilo\hertz} using Betaflight's blackbox logging system, capturing:
\begin{itemize}
  \item Motor RPM (bidirectional DShot telemetry from all four motors)
  \item Attitude (roll, pitch, yaw computed via Mahony AHRS filter from IMU)
  \item GPS velocity and barometric altitude
  \item Battery voltage and current consumption (Hall-effect sensor, accuracy $\pm 2\%$)
\end{itemize}

A representative \SI{10}{\second} cruise segment at approximately \SI{10}{\meter\per\second} forward velocity was extracted for analysis. This segment represents steady-state cruise with manually maintained altitude under calm conditions.

\subsection{Flight Performance Metrics}\label{sec:outdoor-metrics}

Table~\ref{tab:outdoor_power_summary} summarizes the measured electrical power consumption at different flight speeds, demonstrating the progressive power reduction with increasing airspeed.

\begin{table}[htbp]
\centering
\caption{Outdoor flight electrical power measurements showing power reduction with airspeed.}
\label{tab:outdoor_power_summary}
\begin{tabular}{ccccc}
\hline
\textbf{Speed} & \textbf{Avg RPM} & \textbf{Avg I} & \textbf{Avg V} & \textbf{Power} \\
\textbf{[\si{\meter\per\second}]} & & \textbf{[\si{\ampere}]} & \textbf{[\si{\volt}]} & \textbf{[\si{\watt}]} \\
\hline
0 (Hover) & 11957 & 12.0 & 50.0 & 600 \\
10.0      &  9539 &  8.0 & 50.0 & 400 \\
\hline
\multicolumn{5}{l}{\textbf{$\Delta P$ vs Hover:} $-33.3\%$} \\
\hline
\end{tabular}
\end{table}

The data show a clear electrical power reduction during cruise: average current decreased from \SI{12.0}{\ampere} in hover to \SI{8.0}{\ampere} at \SI{10}{\meter\per\second}, corresponding to a 33\% reduction in total electrical power (600~W to 400~W at constant voltage).

\subsection{Time-Series Analysis}\label{sec:outdoor-timeseries}

Figure~\ref{fig:outdoor_timeseries} presents the complete time-series data from the cruise segment, showing $P(t)$ and $\text{RPM}(t)$ together. Both power and motor speed exhibit a smooth decline as airspeed increases, confirming progressive rotor unloading.

\begin{figure}[p]
  \centering
  \input{figures/pgfplots_styles.tex}
  
  % GPS Velocity
  \begin{tikzpicture}
  \begin{axis}[
    tumplot,
    width=0.95\textwidth,
    height=0.15\textheight,
    xlabel={},
    ylabel={GPS Speed [m/s]},
    grid=major,
    title={GPS Velocity},
  ]
  \addplot[TUMBlue, thick] table[x=time_s, y=gps_speed_ms, col sep=comma] {data/outdoor/outdoor_flight_motor_data.csv};
  \end{axis}
  \end{tikzpicture}
  
  \vspace{0.3cm}
  
  % Barometric Altitude
  \begin{tikzpicture}
  \begin{axis}[
    tumplot,
    width=0.95\textwidth,
    height=0.15\textheight,
    xlabel={},
    ylabel={Altitude [m]},
    grid=major,
    title={Barometric Altitude},
  ]
  \addplot[TUMAccentGreen, thick] table[x=time_s, y=altitude_rel_m, col sep=comma] {data/outdoor/outdoor_flight_motor_data.csv};
  \addplot[black, dashed, thick] coordinates {(3, 0) (13, 0)};
  \end{axis}
  \end{tikzpicture}
  
  \vspace{0.3cm}
  
  % Pitch Angle
  \begin{tikzpicture}
  \begin{axis}[
    tumplot,
    width=0.95\textwidth,
    height=0.15\textheight,
    xlabel={},
    ylabel={Pitch Angle [°]},
    grid=major,
    title={Pitch Angle},
  ]
  \addplot[TUMAccentOrange, thick] table[x=time_s, y=pitch_deg, col sep=comma] {data/outdoor/outdoor_flight_motor_data.csv};
  \addplot[black, dashed, thick] coordinates {(3, 0) (13, 0)};
  \end{axis}
  \end{tikzpicture}
  
  \vspace{0.3cm}
  
  % Electrical Power
  \begin{tikzpicture}
  \begin{axis}[
    tumplot,
    width=0.95\textwidth,
    height=0.15\textheight,
    xlabel={},
    ylabel={Power [W]},
    grid=major,
    title={Electrical Power},
  ]
  \addplot[TUMAccentBlue, thick] table[x=time_s, y=power_W, col sep=comma] {data/outdoor/outdoor_flight_motor_data.csv};
  \end{axis}
  \end{tikzpicture}
  
  \vspace{0.3cm}
  
  % Average Motor RPM
  \begin{tikzpicture}
  \begin{axis}[
    tumplot,
    width=0.95\textwidth,
    height=0.15\textheight,
    xlabel={Time [s]},
    ylabel={Motor RPM},
    grid=major,
    title={Average Motor RPM},
  ]
  \addplot[TUMSecondaryBlue2, thick] table[x=time_s, y=motor_avg_rpm, col sep=comma] {data/outdoor/outdoor_flight_motor_data.csv};
  \end{axis}
  \end{tikzpicture}
  
  \caption{Time-series data from outdoor cruise flight at \SI{10}{\meter\per\second}. From top to bottom: GPS velocity maintains steady cruise speed; barometric altitude variations ($\pm$\SI{0.5}{\meter}) demonstrate manual altitude control; pitch angle near \SI{87.5}{\degree} indicates near-level flight; electrical power and average motor RPM both decline smoothly with airspeed, confirming progressive rotor unloading.}
  \label{fig:outdoor_timeseries}
\end{figure}

\subsection{Discussion and Validation}\label{sec:outdoor-discussion}

The outdoor results provide direct evidence of aerodynamic efficiency. Average electrical power decreased from $P_0 = 600$~W in hover to $P_1 = 400$~W at \SI{10}{\meter\per\second}, consistent with the predicted rotor-speed reduction observed in indoor testing. This verifies that the observed rotor unloading translates into measurable power savings.

Measurement uncertainty is small compared to the observed $\Delta P = 33\%$, confirming robustness of the result. The smooth decline in both $P(t)$ and $\text{RPM}(t)$ with increasing airspeed demonstrates that aerodynamic lift progressively offsets rotor thrust, validating the fundamental efficiency advantage of the hybrid platform.

Key findings:
\begin{itemize}
  \item \textbf{Direct power validation}: Electrical power measurements confirm that thrust-based rotor unloading (observed indoors) corresponds to real electrical-power reduction in outdoor flight.
  \item \textbf{Quantified efficiency}: 33\% power reduction at \SI{10}{\meter\per\second} cruise provides a validated baseline for mission planning and endurance estimation.
\end{itemize}

\subsection{Power Model and Theoretical Efficiency Estimation}\label{sec:power-model}

To interpret the measured power reduction, the analytical model proposed by Bauersfeld et al.~\cite{bauersfeld2021range} was used to estimate hover and forward-flight power requirements of a 2.5\,kg quadrotor with 7-inch propellers.
The method provides closed-form estimates for hover, endurance-optimal, and range-optimal flight, validated experimentally for multirotors up to \SI{18}{\meter\per\second}.

\paragraph{Model overview.}
The total mechanical power is the sum of induced and parasitic components:
\begin{equation}
P(V) = P_i(V) + P_p(V),
\end{equation}
where the induced term decreases with speed due to reduced rotor loading, while the parasitic term grows quadratically with airspeed.
Bauersfeld expresses the power ratio relative to hover as
\begin{equation}
\frac{P(V)}{P_h} = k_i(V) + k_p(V),
\end{equation}
with empirical coefficients giving minima at the endurance-optimal speed $V_e$ and range-optimal speed $V_r$.

\paragraph{Baseline hover power.}
For the X-wing's propulsion system
\begin{equation}
m = 2.5\,\text{kg}, \quad r = 0.089\,\text{m}, \quad N_r = 4, \quad \rho = 1.225\,\text{kg\,m}^{-3},
\end{equation}
the induced velocity in hover is
\begin{equation}
v_{i,h} = \sqrt{\frac{mg}{2\rho\pi r^2 N_r}} \approx 10.0\,\text{m\,s}^{-1},
\end{equation}
yielding a theoretical mechanical hover power
\begin{equation}
P_h = \frac{(mg)^{3/2}}{\eta_p\sqrt{2\rho\pi N_r r^2}} \approx 410\,\text{W},
\end{equation}
assuming propeller figure-of-merit $\eta_p = 0.6$.
With motor efficiency $\eta_m = 0.75$, the electrical hover power is roughly 547\,W.

\paragraph{Predicted forward-flight power.}
Using Bauersfeld's experimentally validated ratios
\begin{equation}
\frac{P_e}{P_h} = 0.914, \quad \frac{P_r}{P_h} = 1.092,
\end{equation}
and their corresponding speeds for a vehicle of comparable disk loading,
\begin{equation}
V_e \approx 12.3\,\text{m\,s}^{-1}, \quad V_r \approx 22.1\,\text{m\,s}^{-1},
\end{equation}
the resulting electrical powers are shown in Table~\ref{tab:power_model_comparison}.

\begin{table}[htbp]
\centering
\caption{Theoretical power predictions from Bauersfeld model for X-wing platform.}
\label{tab:power_model_comparison}
\begin{tabular}{lccccc}
\hline
\textbf{Flight Regime} & \textbf{Speed} & \textbf{$P_{\text{mech}}$} & \textbf{$P_{\text{elec}}$} & \textbf{$P/P_h$} \\
& \textbf{[\si{\meter\per\second}]} & \textbf{[\si{\watt}]} & \textbf{[\si{\watt}]} & \textbf{[\%]} \\
\hline
Hover & 0 & 410 & 547 & 100 \\
Endurance-opt. & 12.3 & 375 & 500 & 91 \\
Range-opt. & 22.1 & 448 & 597 & 109 \\
10 (operational) & 10.0 & 382 & 509 & 93 \\
\hline
\end{tabular}
\end{table}

Thus, \SI{10}{\meter\per\second} lies close to the endurance-optimal region, with predicted electrical power about 480--520\,W ($\approx$7\% below hover).

\paragraph{Interpretation.}
For a pure quadrotor, the Bauersfeld model predicts only a modest power decrease ($\approx$7\%) near \SI{10}{\meter\per\second} due to reduced induced power partially offset by parasitic drag.
The measured 33\% reduction therefore cannot be explained by rotor aerodynamics alone and confirms that aerodynamic lift from the fixed wings provides additional load relief.

\paragraph{Conclusion.}
Combining theoretical and experimental data substantiates that the observed efficiency gain stems from passive aerodynamic lift rather than measurement bias or rotor operating point.
The Bauersfeld model provides a defensible lower bound for non-aerodynamic multirotors and quantifies the excess savings attributable to the X-wing configuration.

\section{Discussion}\label{sec:discussion}

The experimental campaign presented in this chapter validates the central hypothesis: an INDI-based controller can maintain stable, accurate flight on a quadrotor with strong, unmodeled aerodynamic lift and drag forces.
The results span four complementary experimental domains—thrust characterization, agility testing, aerodynamic coefficient identification, and outdoor power measurement—providing comprehensive validation of both the control architecture and the platform's efficiency benefits.

\paragraph{Control performance.}
The circular trajectory tests at the Munich flight arena demonstrate robust tracking performance across a wide speed range (\SI{3}{\meter\per\second} to \SI{8}{\meter\per\second}) with RMS position errors below \SI{0.24}{\meter} even at maximum speed.
The platform sustained centripetal accelerations up to \SI{27.8}{\meter\per\second\squared} (approximately 2.8g) without attitude saturation or oscillation, confirming that the INDI controller successfully rejects aerodynamic disturbances across all flight speeds.
Critically, no controller retuning was required between hover and forward flight, demonstrating robustness to large aerodynamic regime changes.
This validates that incremental disturbance compensation can absorb quasi-steady aerodynamic effects without explicit modeling.
The aggressive pitch angles achieved during cruise flight (up to \SI{37.2}{\degree} from vertical) would likely exceed the operational envelope of traditional geometric controllers, which typically lack the disturbance rejection capability necessary to handle large unmodeled aerodynamic moments.

\paragraph{Aerodynamic validation.}
The experimental aerodynamic coefficient identification using force balance analysis confirmed excellent agreement between XFLR5 predictions and flight-test measurements for lift coefficients.
The measured drag coefficients were consistently higher than wing-only simulations, as expected due to system-level drag sources (fuselage, propellers, interference effects).
This validates the computational model for lift estimation while highlighting the importance of considering complete system drag for range and endurance prediction.
The AG25 airfoil demonstrated 43\% higher maximum L/D ratio compared to NACA~0015 (66.4 vs 46.4), confirming the benefit of airfoil optimization for this Reynolds number regime.

\paragraph{Power efficiency.}
Outdoor flight tests with electrical power logging provided direct evidence of aerodynamic efficiency: average electrical power decreased from 600~W in hover to 400~W at \SI{10}{\meter\per\second} cruise, corresponding to a 33\% power reduction.
The smooth, monotonic relationship between airspeed and power confirms progressive rotor unloading due to aerodynamic lift, validating the fundamental efficiency advantage of the hybrid platform.
This measurement bridges the gap between thrust-based indoor analysis and real-world electrical power consumption, providing validated performance metrics for mission planning.

\section{Limitations and Future Work}\label{sec:limitations}

While the experimental results demonstrate successful integration of aerodynamic surfaces with incremental nonlinear control, several limitations suggest directions for future investigation:

\begin{itemize}
  \item \textbf{Limited speed envelope}: Flight tests were restricted to \SI{8}{\meter\per\second} indoors and \SI{10}{\meter\per\second} outdoors due to arena constraints. Higher-speed testing would quantify the limits of quasi-steady aerodynamic assumptions and identify when dynamic effects (e.g., gust response, unsteady flow) require explicit aerodynamic modeling in the controller.
  
  \item \textbf{No controller baseline comparison}: While the platform's aggressive pitch capability and stable flight suggest INDI's superior disturbance rejection compared to traditional geometric controllers, quantitative A/B testing was not performed. Direct comparison with baseline controllers (cascaded PID, geometric, nonlinear dynamic inversion with aerodynamic models) over identical trajectories would isolate the specific contributions of incremental feedback and quantify performance improvements in tracking accuracy, energy efficiency, and control effort.
  
  \item \textbf{Aerodynamic model simplifications}: The XFLR5 simulations neglect propeller--wing aerodynamic interference, which likely affects both lift and drag. Wind-tunnel testing or high-fidelity CFD with rotating propellers would refine force predictions and potentially reveal optimization opportunities for propeller placement or wing positioning.
  
  \item \textbf{Short flight duration}: All indoor tests were limited to \SI{1}{\minute} or less, preventing assessment of thermal effects on propulsion efficiency, battery voltage sag under sustained load, or potential structural issues (e.g., wing vibration, motor bearing wear) during extended operation.
  
  \item \textbf{Manual outdoor flight}: Outdoor tests relied on manual piloting for safety, introducing variability in flight path and attitude control. Future work should implement autonomous outdoor navigation with repeatable trajectories to enable statistical comparison of power consumption across multiple flights.
  
  \item \textbf{Environmental conditions}: Indoor flights eliminate wind disturbances, while outdoor tests were conducted in calm conditions. Systematic testing under varying wind speeds and turbulence levels would quantify the controller's disturbance rejection limits and inform safe operating envelopes.
\end{itemize}

\section{Summary}\label{sec:exp-summary}

This chapter presented comprehensive experimental validation of the X-wing platform across four complementary test campaigns:

\begin{enumerate}
  \item \textbf{Thrust characterization} established the quadratic thrust-to-RPM relationship ($R^2 = 0.997$) required for control allocation and power analysis.
  
  \item \textbf{Circular trajectory tracking} at the Munich flight arena demonstrated sub-decimeter position accuracy (\SI{0.064}{\meter} to \SI{0.184}{\meter} RMS error) across speeds from \SI{3}{\meter\per\second} to \SI{8}{\meter\per\second}, with sustained centripetal accelerations up to 2.8g, validating robust INDI-based control without explicit aerodynamic modeling.
  
  \item \textbf{Aerodynamic coefficient identification} using force balance analysis confirmed XFLR5 lift predictions and quantified system-level drag contributions, with the AG25 airfoil achieving 43\% higher L/D ratio than NACA~0015.
  
  \item \textbf{Outdoor power measurements} provided direct evidence of 33\% electrical power reduction at \SI{10}{\meter\per\second} cruise compared to hover, validating the efficiency benefits of aerodynamic lift.
\end{enumerate}

The experimental evidence demonstrates that incremental nonlinear control can stabilize and accurately control a quadrotor whose aerodynamic surfaces produce lift on the order of the vehicle's weight—without requiring aerodynamic identification or gain scheduling.
The combination of airframe design (optimized airfoil) and control architecture (INDI) delivers measurable efficiency improvements for cruise flight while preserving quadrotor-like agility.
The next chapter concludes with implications for hybrid aerial vehicle design and prospective extensions of this work.
